% Complete sharp description of the first missing level for increasing hosts.
\section{Every first missing word is a hook}
Call $(j,1^{n-j})$ a hook of weight $n$; this includes $1^n$
when $j=1$. Fix an increasing host $X$, and either allow all
parts or restrict them to $1,\ldots,t$, where $t=\max X$.
Let $\rho$ be the first weight of a missing allowed word.

\begin{theorem}[Sharp first-defect support]
Every missing allowed word of weight $\rho$ is a hook.
For the bounded alphabet, the number of missing words at this
rank is at most $\min(\rho-1,t)$. Without the restriction it is
at most $\rho$. The exact number equals the number of failing
hook tests at that rank.
\end{theorem}
\begin{proof}
Put $n=\rho$. The host contains every allowed word of weight
$n-1$, and therefore its last $n-1$ capacities dominate
$\min(1,t),\ldots,\min(n-1,t)$ in the bounded case. Omit
the minima in the unrestricted case.

Let $u$ have weight $n$, length $k$, and fail to be a hook.
Then $k\le n-1$. Every letter is at most $n-k+1$.
If a letter equals this upper bound, all other letters are
ones. In particular the first letter is at most $n-k$:
otherwise $u$ would be a hook. Right-align $u$ with the
last $k$ host positions. In the staircase for weight $n-1$,
the first of these positions has coordinate $n-k$, and all
later ones have coordinates at least $n-k+1$. Thus every
letter fits; in the bounded case also use $u_i\le t$.
This proves that no nonhook can be missing.

There is exactly one hook for each possible first letter.
A hook $(j,1^{n-j})$ embeds precisely when the host has at
least $n-j+1$ positions and the capacity at position
$\ell-n+j$ is at least $j$. These tests give both the
exact count. In the bounded case a one-letter word of an
allowed weight is always present, so a missing hook cannot
have $j=n$. This gives the slightly sharper bounded count.
\end{proof}

\begin{corollary}[The next boundary coefficient]
In the bounded case $t\ge2$, let $a>0$ be the missing-hook
count and let $\varepsilon$ indicate whether $1^\rho$ is
missing. For the boundary kernel $R=1-(1-S_t)F_X$,
\[
                          R_{\rho+1}\ge2a-\varepsilon.
\]
\end{corollary}
\begin{proof}
Every missing hook has a trailing one. The following three
disjoint families are missing words of weight $\rho+1$:
append one to each missing hook; change its last one to two;
and prepend one to every missing hook whose first part is
at least two. The families have sizes $a,a,a-\varepsilon$.
They stay inside the bounded alphabet and remain missing by
downset closure. The first family consists of hooks, the
second ends in two, and the third starts in one and ends in
one; their disjointness and the within-family injections
are immediate. Thus the next missing count is at least
$3a-\varepsilon$. Multiplication by $1-S_t$ subtracts just
the first missing count at this next rank, giving the result.
\end{proof}

Both bounds are sharp. The host $(1,2,\ldots,n-1)$ is
universal through weight $n-1$ and misses all $n$ hooks at
weight $n$. With the alphabet bounded by $n-1$, it misses
exactly the $n-1$ allowed hooks at that rank.

\begin{corollary}[The exact leading finite-boundary coefficient]
Use the maximum-capacity reduction, with first restricted
missing weights $\rho_X,\rho_Y$ and missing-hook counts
$a\le t$, $b\le s$. If
$\mu=\min(\rho_X+s+1,\rho_Y+t+1,\rho_X+\rho_Y)$, then
\[
 \mathcal L_\mu
 =a\mathbf1_{\mu=\rho_X+s+1}
  +b\mathbf1_{\mu=\rho_Y+t+1}
  +ab\mathbf1_{\mu=\rho_X+\rho_Y}
 \le t+s+ts.
\]
The bound is attained with $t=s$, both arms $(1,2,\ldots,t)$,
and any strict peak $M>t$.
\end{corollary}
\begin{proof}
The leading coefficients of the two boundary kernels are
the first missing counts, because the unrestricted comparison
series have constant coefficient one. Apply the exact loss
identity; its negative term begins one rank later. For the
displayed staircase arms, $\rho_X=\rho_Y=t+1$ and $a=b=t$,
so all three leading ranks coincide and give $t^2+2t$.
\end{proof}
